\section{Ley de Control Nominal}

La cinemática del robot diferencial en coordenadas cartesianas es:
\begin{equation}
  \dot{x} = v\cos\theta, \quad \dot{y} = v\sin\theta, \quad \dot{\theta} = \omega,
  \label{eq:cinematica}
\end{equation}
donde $(x, y, \theta)$ es la posición y orientación del robot, y $(v, \omega)$ son las velocidades lineal y angular.

Dado un punto objetivo $(x_r, y_r)$, el error de posición en coordenadas polares es:
\begin{equation}
  a = \sqrt{(x_r - x)^2 + (y_r - y)^2}, \quad
  \alpha = \text{atan2}(y_r - y,\; x_r - x) - \theta,
  \label{eq:error_polar}
\end{equation}
donde $a$ es la distancia al objetivo y $\alpha$ es el error de orientación normalizado a $(-\pi, \pi]$.

La ley de control nominal propuesta genera velocidades de referencia:
\begin{equation}
  v_{\text{ref}} = k_1\, a\cos\alpha, \quad
  w_{\text{ref}} = k_2\,\alpha + k_1\sin\alpha\cos\alpha,
  \label{eq:ley_control}
\end{equation}
con ganancias $k_1 = 0.5$ y $k_2 = 1.0$, y saturación en $[-v_{\max},\, v_{\max}]$ y $[-\omega_{\max},\, \omega_{\max}]$ respectivamente. Esta ley garantiza convergencia exponencial al objetivo en ausencia de obstáculos~\cite{ramirez2001collision}.
